\section{Introduction}


\begin{figure}[ht]
    \begin{center}
  \includegraphics[width=1.0\linewidth]{figures/openfig.pdf}
    \end{center}
       \caption{Comparison of  LiDAR scenes generated by different methods on the nuScenes val split. \sysname\ is the first work to achieve sequential LiDAR scene generation with highly controllable scene manipulation abilities, including foreground control, background control and object-fidelity enhancement.}
\label{fig:fig_openfig}
 \end{figure}

Data is a foundational element driving artificial intelligence advances.  Within autonomous driving research, high-quality data is particularly crucial due to: i) the inherent data-intensive requirements of deep learning models, and ii) the necessity of capturing corner cases — rare driving behaviours and uncommon road environments — which are essential for developing safety-critical systems. However, collecting and annotating diverse multi-modal datasets (e.g., camera and LiDAR) remains time-consuming and resource-intensive. While recent generative models have demonstrated promising capabilities for synthesizing visual data, LiDAR scene generation—despite its critical role in providing geometric awareness — remains comparatively underdeveloped.  In this work, we aim to advance existing LiDAR scene generation techniques to better address real-world autonomous driving requirements.



% objects control - road sketch control
To synthesize realistic LiDAR data that accurately captures diverse real-world traffic scenarios, a LiDAR scene generation method should suppport flexible customization of road layouts and dynamic object placements. Recent studies, such as LiDARGen~\cite{zyrianov2022learning}, UltraLiDAR~\cite{xiong2023ultralidar}, R2DM~\cite{nakashima2024lidar}, and RangeLDM~\cite{hu2025rangeldm}, have made considerable advancements in producing realistic LiDAR data.  However, these techniques predominantly generate data in an unconditional manner, lacking the capability to manipulate specific scene elements.


% background control - scene caption control
The quality of background LiDAR point clouds is as critical as that of foreground objects to ensure realism.
To this end, Text2LiDAR~\cite{wu2024text2lidar} utilizes textual descriptions as conditioning input. Nonetheless, this method is confined to coarse descriptions encompassing weather conditions, time of day, and object names. The omission of detailed background information, such as specific descriptions of trees or buildings, compromises the realism of the generated LiDAR data. 

% sequential scene generation
In addition, current 3D LiDAR scene generation methods prove insufficient when it comes to capturing the dynamic behaviours of objects. To mitigate this issue, LidarDM \cite{zyrianov2024lidardm} proposes to generate LiDAR sequences by separately modeling static scenes and dynamic objects. However, it lacks background control, and its two-stage compositional strategy potentially compromises the realism and coherence of the resulting point cloud distributions.


In summary, current LiDAR scene generation methods are notably deficient when it comes to integrating all critical capbalities: i) sequential scene generation with detailed control over both ii) foreground and iii) background components.  
To bridge this gap, we propose \sysname, an end-to-end 4D LiDAR scene generation pipeline that facilitates the generation of sequential LiDAR scenes with comprehensive scene manipulation capabilities. 
\sysname\ is distinguished by two principal features: i) the incorporation of multimodal conditions, including scene captions, road sketches and object priors, and ii) a meticulously designed equirectangular spatial-temporal noise prediction model, \modelname, which ensures spatial and temporal consistency throughout denoising processes. 

\cref{fig:fig_openfig} displays LiDAR scenes generated by different methods on the nuScenes val split. It is evident that RangeLDM \cite{hu2025rangeldm} and Text2LiDAR \cite{wu2024text2lidar} are unable to generate the vehicle accurately, and backgrounds do not align with those observed in the Ground Truth(GT) scene. In contrast, the vehicles' positions and structures, and the background generated by \sysname\ closely align with those in the GT scene. Furthermore, it is noteworthy that \sysname\ is capable of generating a sequence of LiDAR scenes that maintain temporal consistency, whereas RangeLDM \cite{hu2025rangeldm} and Text2LiDAR \cite{wu2024text2lidar} can only create individual LiDAR scenes in isolation.

To summarize, our specific contributions are as follows:
\begin{itemize}
    \item This is the first work to achieve precise control over foreground objects, including manipulation of their positions and sizes, and fine-grained control over background elements in LiDAR scene generation.  
    \item We propose a novel equirectangular spatio-temporal diffusion model, \modelname,  which achieves end-to-end \textit{sequential} LiDAR scene generation while ensuring consistency over foreground and background elements.
    \item We demonstrate the effectiveness of our proposed \sysname\ on the KITTI and nuScenes datasets, where it outperforms the current SOTA methods.
\end{itemize}

